
%% ===================================================================
%% LP-134 canonical naming (2025-12-15):
%% Per LP-134 (Lux Chain Topology), M-Chain has been split into:
%%   M-Chain — MPC ceremonies for bridge custody of EXTERNAL wallets
%%             (CGGMP21, FROST, Pulsar-general)
%%   F-Chain — FHE compute + TFHE bootstrap-key generation (encrypted EVM)
%% M-Chain is REMOVED entirely; teleport is B-Chain (bridgevm, LP-6000). No teleportvm.
%% Where this paper says "M-Chain MPC" / "M-Chain threshold governance" /
%% "M-Chain TFHE" / "M-Chain custody", read as M-Chain (MPC) or F-Chain (FHE).
%% ===================================================================

\section{Agent Capabilities}
\label{sec:capabilities}

The capability model governs what an agent is permitted to do. Capabilities are fine-grained, composable, and enforced at the protocol level.

\begin{definition}[Capability Token]
A capability token $\kappa = (r, s, t_0, t_1, \sigma)$ comprises:
\begin{itemize}
    \item $r \in \{\texttt{read}, \texttt{write}, \texttt{sign}, \texttt{spend}, \texttt{delegate}, \texttt{invoke}\}$: the right,
    \item $s$: the scope (a resource identifier or pattern),
    \item $[t_0, t_1]$: the validity window,
    \item $\sigma$: the issuer's signature over $(r, s, t_0, t_1)$.
\end{itemize}
\end{definition}

\subsection{Capability Classes}

We define six capability classes:

\begin{table}[H]
\centering
\caption{Agent capability classes and their scopes.}
\label{tab:capabilities}
\begin{tabular}{llp{6cm}}
\toprule
\textbf{Right} & \textbf{Example Scope} & \textbf{Description} \\
\midrule
\texttt{read} & \texttt{vault://agent/\{id\}/*} & Read data from a specified vault path \\
\texttt{write} & \texttt{vault://agent/\{id\}/memory/*} & Write to the agent's memory namespace \\
\texttt{sign} & \texttt{chain://c-chain/tx/*} & Sign transactions on a specific chain \\
\texttt{spend} & \texttt{budget://\{id\}/max:100LUX} & Spend up to a specified amount \\
\texttt{delegate} & \texttt{cap://\{id\}/depth:2} & Delegate capabilities up to a depth \\
\texttt{invoke} & \texttt{mcp://tool/web-search} & Invoke a specific MCP tool \\
\bottomrule
\end{tabular}
\end{table}

\subsection{Threshold Approval}

High-value actions require threshold approval rather than unilateral agent execution. A threshold policy $\pi = (n, k, \mathcal{A})$ specifies that $k$ out of $n$ authorized parties in set $\mathcal{A}$ must co-sign the action. Threshold signing is performed via the M-Chain's distributed key generation (DKG) and threshold signature scheme~\cite{gennaro2020threshold}:

\begin{equation}\label{eq:threshold}
    \sigma_{\text{threshold}} = \text{TSS.Combine}\big(\{\sigma_i\}_{i \in S}, |S| \geq k\big)
\end{equation}

This prevents a compromised agent from unilaterally executing catastrophic actions (e.g., draining its entire budget, signing high-value transactions, delegating all capabilities).

\subsection{Time-Bounded Grants}

All capabilities expire. There are no permanent grants. The maximum validity window $t_1 - t_0$ is bounded by a protocol parameter $\Delta_{\max}$ (default: 30 days). Renewal requires explicit re-issuance by the granting principal, forcing periodic review of agent permissions.

